% ============================================================
% CAPITOLUL 2 — EPIDEMIOLOGIE COMPUTAȚIONALĂ ȘI MODELAREA
%               MATEMATICĂ A BOLILOR INFECȚIOASE
% ============================================================

\chapter{Epidemiologie Computațională și Modelarea Matematică a Bolilor Infecțioase}
\label{ch:epidemiologie}

\section{Fundamente de epidemiologie}
\label{sec:fundamente_epi}

\subsection{Definiții și concepte fundamentale}
\label{subsec:definitii_epi}

Epidemiologia este știința care studiază distribuția, determinanții și controlul
stărilor sau evenimentelor legate de sănătate în populații definite
\cite{andersonmay1991}. În sens mai specific, \emph{epidemiologia infecțioasă}
se ocupă cu dinamica transmiterii bolilor cauzate de agenți patogeni biologici
(virusuri, bacterii, paraziți) în populații umane sau animale.

\begin{definitie}[Epidemic, pandemic, endemic]
O boală infecțioasă se numește \textbf{endemică} dacă persistă la nivel constant
într-o regiune geografică definită. Se numește \textbf{epidemie} dacă numărul de
cazuri depășește clar nivelul endemic așteptat, și \textbf{pandemie}
dacă afectează mai multe continente sau o mare parte a globului.
\end{definitie}

\begin{definitie}[Incidență și prevalență]
\textbf{Incidența} reprezintă numărul de cazuri noi ale unei boli apărute într-o
populație de risc, într-un interval de timp specificat. Se exprimă ca rată:
\begin{equation}
    \text{Incidență} = \frac{\text{cazuri noi în intervalul } [t_1, t_2]}
                            {\text{populație la risc} \times (t_2 - t_1)}.
\end{equation}
\textbf{Prevalența} reprezintă proporția indivizilor afectați la un moment dat:
\begin{equation}
    \text{Prevalență} = \frac{\text{cazuri existente la momentul } t}
                             {\text{populație totală la momentul } t}.
\end{equation}
\end{definitie}

\subsection{Parametrii fundamentali ai transmiterii}
\label{subsec:parametri}

\begin{definitie}[Intervalul serial și intervalul de generație]
\textbf{Intervalul serial} este durata medie dintre apariția simptomelor la cazul
primar și la cazul secundar generat de acesta. \textbf{Intervalul de generație}
este durata medie dintre infecția cazului primar și infecția unui caz secundar.
\end{definitie}

\begin{definitie}[Perioada de incubație]
\textbf{Perioada de incubație} este intervalul de timp dintre momentul infectării
și apariția primelor simptome. Este distinctă de \emph{perioada de latență},
care reprezintă intervalul dintre infecție și debutul infecțiozității.
\end{definitie}

\begin{definitie}[Numărul de reproducere de bază $\Rzero$]
\label{def:R0}
Numărul de reproducere de bază $\Rzero$ este definit ca numărul mediu de infecții
secundare produse de un singur individ infectat introdus într-o populație
complet susceptibilă, pe durata perioadei sale infecțioase \cite{diekmann1990}.
Este cel mai important parametru al epidemiologiei matematice, deoarece determină
potențialul epidemic al unui agent patogen.
\end{definitie}

O consecință imediată a \cref{def:R0} este existența unui prag epidemic. Dacă
$\Rzero < 1$, fiecare caz infectat produce în medie mai puțin de un caz nou, iar
epidemia se stinge spontan. Dacă $\Rzero > 1$, numărul de cazuri crește în faza
inițială, de regulă cu profil aproximativ exponențial. Cazul $\Rzero = 1$ separă
cele două regimuri și corespunde unei stări de echilibru instabil.

Numărul de reproducere \emph{efectiv} (sau în timp real) $\mathcal{R}(t)$ generalizează
$\Rzero$ pentru o populație parțial imunizată:
\begin{equation}
    \mathcal{R}(t) = \Rzero \cdot \frac{S(t)}{N},
\end{equation}
unde $S(t)$ este numărul de indivizi susceptibili la momentul $t$ și $N$ este
mărimea populației totale.

% ============================================================
\section{Modele compartimentale în epidemiologie}
\label{sec:modele_compartimentale}
% ============================================================

\subsection{Cadrul general}
\label{subsec:cadru_general}

Modelele compartimentale reprezintă clasa fundamentală de modele matematice în
epidemiologie \cite{hethcote2000}. Principiul lor constă în împărțirea populației
în \emph{compartimente} care descriu stări epidemiologice distincte, iar tranziția
dintre compartimente este guvernată de un sistem de ecuații diferențiale ordinare
(EDO).

Fie $N$ mărimea totală a populației (presupusă constantă). Fiecare compartiment
$X_i(t)$ reprezintă numărul de indivizi în starea $i$ la momentul $t$, cu
$\sum_i X_i(t) = N$. Dinamic\-a tranzițiilor este descrisă de sistemul:
\begin{equation}
    \diff{\mathbf{x}}{t} = \mathbf{f}(\mathbf{x}, t, \boldsymbol{\theta}),
    \quad \mathbf{x}(0) = \mathbf{x}_0,
\end{equation}
unde $\mathbf{x} = (X_1, \ldots, X_n)^\intercal$, $\boldsymbol{\theta}$ este
vectorul parametrilor și $\mathbf{f}$ descrie fluxurile nete dintre compartimente.

\subsection{Modelul SIS}
\label{subsec:model_sis}

Modelul \textbf{SIS} (Susceptibili-Infectați-Susceptibili) descrie infecții în
care recuperarea nu conferă imunitate de durată, permițând re-infectarea:
\begin{equation}
\label{eq:sis}
\begin{cases}
    \diff{S}{t} = -\dfrac{\beta S I}{N} + \gamma I, \\[8pt]
    \diff{I}{t} = \phantom{-}\dfrac{\beta S I}{N} - \gamma I.
\end{cases}
\end{equation}
Parametrul $\gamma > 0$ reprezintă rata de recuperare ($1/\gamma$ este durata
medie a perioadei infecțioase). Pentru modelul SIS, $\Rzero = \beta/\gamma$.
Dacă $\Rzero \leq 1$, echilibrul fără boală $I^* = 0$ este stabil global; dacă
$\Rzero > 1$, sistemul admite un echilibru endemic stabil
$I^* = N(1 - 1/\Rzero) > 0$.

\subsection{Modelul SEIR}
\label{subsec:model_seir}

Modelul \textbf{SEIR} introduce compartimentul $E$ (Expuși / în perioadă de latență),
care ține cont de faptul că indivizii infectați nu sunt imediat infecțioși:
\begin{equation}
\label{eq:seir}
\begin{cases}
    \diff{S}{t} = -\dfrac{\beta S I}{N}, \\[8pt]
    \diff{E}{t} = \phantom{-}\dfrac{\beta S I}{N} - \sigma E, \\[8pt]
    \diff{I}{t} = \phantom{-}\sigma E - \gamma I, \\[8pt]
    \diff{R}{t} = \phantom{-}\gamma I,
\end{cases}
\end{equation}
unde $\sigma = 1/T_E$ este rata de tranziție din starea expusă în starea infecțioasă
($T_E$ = durata medie de incubație). Numărul de reproducere de bază pentru modelul
SEIR este $\Rzero = \beta/\gamma$, identic cu cel al modelului SIR.

\subsection{Modelul SIRS}
\label{subsec:model_sirs}

Modelul \textbf{SIRS} generalizează modelul SIR prin imunitate temporară (indivizii
recuperați pot deveni din nou susceptibili după un interval de timp):
\begin{equation}
\label{eq:sirs}
\begin{cases}
    \diff{S}{t} = -\dfrac{\beta S I}{N} + \xi R, \\[8pt]
    \diff{I}{t} = \phantom{-}\dfrac{\beta S I}{N} - \gamma I, \\[8pt]
    \diff{R}{t} = \phantom{-}\gamma I - \xi R,
\end{cases}
\end{equation}
unde $\xi = 1/T_R$ este rata de pierdere a imunității ($T_R$ = durata medie a
imunității). Dacă $\Rzero > 1$, modelul SIRS admite un echilibru endemic stabil.

% ============================================================
\section{Modelul SIDARTHE}
\label{sec:model_sidarthe}
% ============================================================

Modelul SIDARTHE, propus pentru analiza epidemiei COVID-19
\cite{giordano2020sidarthe}, separă două dimensiuni care lipsesc din SIR:
statutul de diagnosticare și severitatea clinică. Într-un sistem de raportare
real, persoanele infectate dar nedepistate, persoanele diagnosticate, cazurile
simptomatice și cazurile severe nu au aceeași contribuție la transmitere și nu
sunt observate cu aceeași probabilitate. Sistemul implementat în această lucrare
folosește varianta \textbf{SIDARTHE}, care include explicit compartimentul de
decese.

Compartimentele modelului sunt $S$ (susceptibili), $I$ (infectați nedepistați),
$D$ (diagnosticați asimptomatici), $A$ (simptomatici nedepistați), $R$
(simptomatici diagnosticați), $T$ (cazuri severe), $H$ (vindecați) și $E$
(extinși / decedați). Compartimentul $E$ este alimentat din $T$ prin rata de
mortalitate $\tau$, ceea ce permite analiza separată a vindecării și a deceselor.

\begin{align}
\diff{S}{t} &= - S \cdot \frac{\alpha I + \beta D + \gamma A + \delta R}{N}, \\
\diff{I}{t} &= S \cdot \frac{\alpha I + \beta D + \gamma A + \delta R}{N}
              - (\varepsilon + \zeta + \lambda) I, \\
\diff{D}{t} &= \varepsilon I - (\eta + \rho)D, \\
\diff{A}{t} &= \zeta I - (\theta + \mu + \kappa)A, \\
\diff{R}{t} &= \eta D + \theta A - (\nu + \xi)R, \\
\diff{T}{t} &= \mu A + \nu R - (\sigma + \tau)T, \\
\diff{H}{t} &= \lambda I + \rho D + \kappa A + \xi R + \sigma T, \\
\diff{E}{t} &= \tau T.
\label{eq:sidarthe}
\end{align}

Parametrii $\alpha$, $\beta$, $\gamma$ și $\delta$ controlează contribuția la
transmitere a compartimentelor infecțioase. Parametrii $\varepsilon$, $\zeta$ și
$\lambda$ descriu ieșirile din starea infectată nedepistată: diagnosticare,
apariția simptomelor și vindecare fără diagnostic. Restul parametrilor descriu
tranzițiile către forme diagnosticate, severe sau vindecate. Spre deosebire de
SIR, unde aproape toată interpretarea se concentrează în raportul
$\beta/\gamma$, SIDARTHE distribuie dinamica pe un număr mai mare de rate; de
aceea, analiza numerică și sensibilitatea la parametri devin mai importante decât
obținerea unei formule scalare simple.

Conservarea populației se verifică prin însumarea ecuațiilor din
\eqref{eq:sidarthe}. Toți termenii interni se anulează, iar relația devine
\begin{equation}
\diff{}{t}(S+I+D+A+R+T+H+E)=0.
\end{equation}
Prin urmare, pentru condiții inițiale nenegative, suma compartimentelor rămâne
egală cu populația totală $N$. Pozitivitatea soluției se
obține prin același argument folosit la modelele compartimentale clasice: pe
frontiera fiecărui compartiment, derivata nu poate împinge soluția în regiunea
negativă, deoarece termenii de ieșire dispar când variabila respectivă este zero.

% ============================================================
\section{Modelul SIR ca reper secundar}
\label{sec:model_sir}
% ============================================================

\subsection{Ipoteze și derivarea ecuațiilor din principii prime}
\label{subsec:sir_derivare}

Modelul SIR, formulat de Kermack și McKendrick \cite{kermack1927}, se bazează
pe o serie de ipoteze simplificatoare care trebuie înțelese explicit:

\begin{enumerate}
    \item \textbf{Populație închisă și omogenă:} Nu există nașteri, decese
    naturale sau migrații; $N = \text{const}$. Fiecare individ are aceeași
    probabilitate de a intra în contact cu orice alt individ (amestecare omogenă).
    \item \textbf{Transmitere prin acțiune de masă:} Rata de contact eficient
    (producătoare de infecție) este proporțională cu produsul $S \cdot I$.
    \item \textbf{Imunitate permanentă:} Un individ recuperat nu poate fi
    re-infectat; compartimentul $R$ este absorbant.
    \item \textbf{Durata infecțioasă exponențial distribuită:} Rata de recuperare
    $\gamma$ este constantă, iar durata medie a perioadei infecțioase este $1/\gamma$.
    \item \textbf{Fără mortalitate indusă de boală:} Sau, echivalent, mortalitatea
    specifică bolii este neglijabilă față de dinamica transmiterii.
\end{enumerate}

\textbf{Derivarea termenului de infecție.} Considerăm că în unitatea de timp,
un individ susceptibil are în medie $c$ contacte cu alți indivizi aleși uniform
din populație. Probabilitatea ca un contact arbitrar să fie cu un individ infectat
este $I/N$, iar probabilitatea de transmitere la contact eficient este $p$. Astfel,
rata de infecție a unui susceptibil este $c \cdot p \cdot I/N = \beta I / N$, unde
$\beta = cp$ este parametrul de transmitere. Fluxul total de noi infecții pe unitate
de timp este $(\beta I / N) \cdot S$.

Sistemul de ecuații diferențiale ordinare al modelului SIR este:
\begin{equation}
\label{eq:sir}
\boxed{
\begin{aligned}
    \diff{S}{t} &= -\frac{\beta S I}{N}, \\[6pt]
    \diff{I}{t} &= \frac{\beta S I}{N} - \gamma I, \\[6pt]
    \diff{R}{t} &= \gamma I,
\end{aligned}
}
\end{equation}
cu condiții inițiale $S(0) = S_0 > 0$, $I(0) = I_0 > 0$, $R(0) = R_0 \geq 0$
și $S_0 + I_0 + R_0 = N$.

Diagrama de flux a compartimentelor este prezentată în
\cref{fig:sir_diagram}.

\begin{figure}[htbp]
\centering
\resizebox{\columnwidth}{!}{%
\begin{tikzpicture}[
    box/.style = {rectangle, draw, thick, minimum width=2cm, minimum height=1.2cm,
                  fill=blue!10, font=\large\bfseries},
    arr/.style = {-{Stealth[length=3mm]}, thick},
    lbl/.style = {font=\normalsize, midway, above}
]
    \node[box] (S) {$S$};
    \node[box, right=3cm of S] (I) {$I$};
    \node[box, right=3cm of I] (R) {$R$};

    \draw[arr] (S) -- (I) node[lbl] {$\beta SI/N$};
    \draw[arr] (I) -- (R) node[lbl] {$\gamma I$};
\end{tikzpicture}
}
\caption{Diagrama de flux a compartimentelor modelului SIR.
Fluxul $\beta SI/N$ reprezintă rata de noi infecții,
iar $\gamma I$ reprezintă rata de recuperare.}
\label{fig:sir_diagram}
\end{figure}

\subsection{Proprietăți matematice fundamentale}
\label{subsec:sir_proprietati}

Înainte de a analiza comportamentul calitativ al soluțiilor, stabilim proprietățile
matematice esențiale: existența și unicitatea soluției, conservarea populației totale
și invarianța domeniului pozitiv.

\begin{teorema}[Existență și unicitate globală]
\label{thm:existenta_unicitate}
Fie $\beta, \gamma > 0$, $N > 0$ și $(S_0, I_0, R_0) \in \R^3$ cu
$S_0, I_0 > 0$, $R_0 \geq 0$. Atunci sistemul \eqref{eq:sir} admite o
unică soluție globală $(S, I, R) \in C^1([0, \infty); \R^3)$.
\end{teorema}

\begin{proof}
Definim $\mathbf{x} = (S, I, R)^\intercal$ și
$\mathbf{f}(\mathbf{x}) = (-\beta SI/N,\; \beta SI/N - \gamma I,\; \gamma I)^\intercal$.
Funcția $\mathbf{f} : \R^3 \to \R^3$ este de clasă $C^\infty$ (polinomul de
gradul doi în $S$ și $I$ este neted), în particular continuu diferențiabilă.
Prin \emph{teorema Picard-Lindelöf} (Cauchy-Lipschitz), există o soluție locală
unică pe un interval $[0, T_{\max})$. \\
\indent Existența \emph{globală} ($T_{\max} = \infty$) se stabilește prin mărginirea
soluției. Deoarece $S(t), I(t), R(t) \geq 0$ (Teorema~\ref{thm:pozitivitate}) și
$S + I + R = N$ (Teorema~\ref{thm:conservare}), rezultă că $\|\mathbf{x}(t)\| \leq N\sqrt{3}$
pentru orice $t \in [0, T_{\max})$. Soluția mărginită nu poate prezenta blow-up în
timp finit, deci $T_{\max} = \infty$.
\end{proof}

\begin{teorema}[Conservarea populației totale]
\label{thm:conservare}
Pentru orice soluție a sistemului \eqref{eq:sir},
\begin{equation}
    S(t) + I(t) + R(t) = N \quad \forall\, t \geq 0.
\end{equation}
\end{teorema}

\begin{proof}
Calculăm derivata sumei:
\begin{align*}
    \diff{}{t}(S + I + R) &= \diff{S}{t} + \diff{I}{t} + \diff{R}{t} \\
    &= -\frac{\beta S I}{N} + \left(\frac{\beta S I}{N} - \gamma I\right) + \gamma I \\
    &= 0.
\end{align*}
Prin integrare, $S(t) + I(t) + R(t) = S_0 + I_0 + R_0 = N$.
\end{proof}

\begin{teorema}[Pozitivitate și mărginire]
\label{thm:pozitivitate}
Dacă $S_0 > 0$, $I_0 > 0$, $R_0 \geq 0$, atunci:
\begin{equation}
    S(t) > 0, \quad I(t) > 0, \quad R(t) \geq 0, \quad \forall\, t > 0.
\end{equation}
În particular, $S(t), I(t), R(t) \in [0, N]$ pentru orice $t \geq 0$.
\end{teorema}

\begin{proof}
\textbf{(a) Pozitivitatea lui $S$:} Din ecuația $dS/dt = -(\beta I/N) S$,
aceasta este o ecuație diferențială liniară omogenă în $S$ cu
coeficient $-\beta I(t)/N$. Soluția este:
\begin{equation}
\label{eq:S_formula}
    S(t) = S_0 \cdot \exp\!\left(-\frac{\beta}{N}\int_0^t I(s)\,ds\right).
\end{equation}
Deoarece integrala este finită și $S_0 > 0$, rezultă $S(t) > 0$ pentru orice $t \geq 0$
(presupunând $I(s) \geq 0$, demonstrat mai jos).

\textbf{(b) Pozitivitatea lui $I$:} Rescrierea ecuației a doua ca
$dI/dt = (\beta S(t)/N - \gamma)\, I$ arată că aceasta este liniară în $I$:
\begin{equation}
    I(t) = I_0 \cdot \exp\!\left(\int_0^t \left[\frac{\beta S(\tau)}{N} - \gamma\right]d\tau\right) > 0,
\end{equation}
deoarece $I_0 > 0$ și exponențiala este strict pozitivă.

\textbf{(c) Non-negativitatea lui $R$:} Din $dR/dt = \gamma I(t) > 0$ și $R_0 \geq 0$,
funcția $R$ este strict crescătoare, deci $R(t) > R_0 \geq 0$.

Mărginirea urmează din Teorema~\ref{thm:conservare}: $S(t) \leq N$, $I(t) \leq N$,
$R(t) \leq N$.
\end{proof}

\begin{observatie}
Teorema~\ref{thm:pozitivitate} implică că domeniul pozitiv
\begin{equation}
    \Omega = \{(S, I, R) \in \R^3_{\geq 0} : S + I + R = N\}
\end{equation}
este \emph{invariant pozitiv} față de fluxul sistemului \eqref{eq:sir}. Prin urmare,
este suficient să studiem comportamentul soluțiilor pe $\Omega$, care este un simplex
în $\R^3$.
\end{observatie}

\begin{teorema}[Monotonia compartimentelor]
\label{thm:monotonie}
De-a lungul oricărei soluții a sistemului \eqref{eq:sir}:
\begin{enumerate}[(i)]
    \item $S(t)$ este strict descrescătoare;
    \item $R(t)$ este strict crescătoare;
    \item $I(t)$ atinge un singur maxim în $t^* = \inf\{t > 0 : S(t) = \gamma N/\beta\}$
    (dacă $S_0 > \gamma N/\beta$); altfel este descrescătoare.
\end{enumerate}
\end{teorema}

\begin{proof}
\textit{(i)} $dS/dt = -\beta SI/N < 0$ deoarece $S > 0$, $I > 0$.
\textit{(ii)} $dR/dt = \gamma I > 0$.
\textit{(iii)} $dI/dt = (\beta S/N - \gamma)I > 0 \Leftrightarrow S > \gamma N/\beta$.
Deoarece $S$ este strict descrescătoare, există cel mult un moment $t^*$ când
$S(t^*) = \gamma N/\beta$, care este maximul lui $I$.
\end{proof}

\subsection{Analiza echilibrelor și stabilitate}
\label{subsec:sir_stabilitate}

Punctele de echilibru ale sistemului \eqref{eq:sir} sunt obținute din condiția
$\mathbf{f}(\mathbf{x}^*) = \mathbf{0}$. Ecuațiile conduc la:
\begin{equation}
    -\frac{\beta S^* I^*}{N} = 0, \quad
    \frac{\beta S^* I^*}{N} - \gamma I^* = 0, \quad
    \gamma I^* = 0.
\end{equation}
Din a treia ecuație, $I^* = 0$. Substituind, obținem că orice punct de forma
$(S^*, 0, N - S^*)$ cu $S^* \in [0, N]$ este echilibru. Mulțimea echilibrelor este:
\begin{equation}
    \mathcal{E} = \{(S^*, 0, N - S^*) : S^* \in [0, N]\} \subset \Omega.
\end{equation}

\begin{observatie}
Spre deosebire de modelele SIS și SIRS, modelul SIR nu admite echilibru endemic
cu $I^* > 0$. Aceasta reflectă faptul că, pe termen lung, epidemia se stinge
întotdeauna (deși poate infecta o fracție mare din populație înainte de stingere).
\end{observatie}

\textbf{Analiza stabilității locale.} Deoarece $R = N - S - I$, sistemul se reduce
la două ecuații (eliminând $R$):
\begin{equation}
\label{eq:sir2d}
\begin{cases}
    \diff{S}{t} = -\dfrac{\beta SI}{N}, \\[8pt]
    \diff{I}{t} = \dfrac{\beta SI}{N} - \gamma I.
\end{cases}
\end{equation}
Jacobianul sistemului \eqref{eq:sir2d} evaluat la un echilibru $(S^*, 0)$ este:
\begin{equation}
\label{eq:jacobian}
    J(S^*, 0) = \begin{pmatrix}
        0 & -\beta S^*/N \\
        0 & \beta S^*/N - \gamma
    \end{pmatrix}.
\end{equation}
Valorile proprii sunt $\lambda_1 = 0$ și $\lambda_2 = \beta S^*/N - \gamma$.

\begin{teorema}[Stabilitatea echilibrului fără boală]
\label{thm:stabilitate_dfe}
Echilibrul fără boală $(S^*, 0)$ este stabil nelinear dacă
$\beta S^*/N < \gamma$, adică dacă $S^*/N < 1/\Rzero$, și este instabil dacă
$\beta S^*/N > \gamma$.
\end{teorema}

\begin{proof}
Valorea proprie $\lambda_2 = \beta S^*/N - \gamma$ determină stabilitatea pe
direcția $I$. Deoarece $\lambda_1 = 0$, analiza liniară este degenerată și
trebuie completată cu argumente neliniare. Totuși, din Teorema~\ref{thm:monotonie},
$I(t)$ este descrescătoare dacă $S(t) < \gamma N/\beta$ pentru tot $t > 0$, ceea
ce se întâmplă când $S_0 < \gamma N/\beta$, implicând stingerea epidemiei.
\end{proof}

\subsection{Numărul de reproducere de bază \texorpdfstring{$\Rzero$}{R0}}
\label{subsec:R0}

\subsubsection{Definiție prin matricea generației următoare}

Metoda \emph{matricei generației următoare} (next-generation matrix), formalizată
de Diekmann, Heesterbeek și Metz \cite{diekmann1990} și extinsă de
van den Driessche și Watmough \cite{vandendriesschwatmough2002},
furnizează o procedură sistematică pentru calculul $\Rzero$ în modele
compartimentale de complexitate arbitrară.

Fie $\mathbf{x}_I$ vectorul compartimentelor infectate. Scriem:
\begin{equation}
    \diff{\mathbf{x}_I}{t} = \mathcal{F}(\mathbf{x}) - \mathcal{V}(\mathbf{x}),
\end{equation}
unde $\mathcal{F}_i$ reprezintă rata de apariție a noilor infecții în
compartimentul $i$, iar $\mathcal{V}_i$ reprezintă rata netă de ieșire
din compartimentul $i$ prin toate celelalte procese.

Matricele Jacobiene evaluate la echilibrul fără boală $\mathbf{x}_0$ sunt:
\begin{equation}
    F = \left[\pdiff{\mathcal{F}_i}{x_j}\right]_{\mathbf{x}_0}, \quad
    V = \left[\pdiff{\mathcal{V}_i}{x_j}\right]_{\mathbf{x}_0}.
\end{equation}
Matricea generației următoare este $K = FV^{-1}$, iar:
\begin{equation}
\label{eq:R0_def}
    \Rzero = \rho(K),
\end{equation}
unde $\rho(\cdot)$ este raza spectrală (cea mai mare valoare proprie în modul).

\subsubsection{Calculul $\Rzero$ pentru modelul SIR}

Pentru modelul SIR, singurul compartiment infectat este $\{I\}$.
La echilibrul fără boală $(S, I, R) = (N, 0, 0)$:
\begin{equation}
    \mathcal{F}(I) = \frac{\beta S I}{N} \implies
    F = \left[\pdiff{\mathcal{F}}{I}\right]_{S=N} = \frac{\beta N}{N} = \beta.
\end{equation}
\begin{equation}
    \mathcal{V}(I) = \gamma I \implies V = \gamma.
\end{equation}
Matricea $K$ este scalarul:
\begin{equation}
    K = F V^{-1} = \frac{\beta}{\gamma}.
\end{equation}
Prin urmare:
\begin{equation}
\label{eq:R0_sir}
    \boxed{\Rzero = \frac{\beta}{\gamma}.}
\end{equation}

Interpretarea biologică este directă: $1/\gamma$ este durata medie a perioadei
infecțioase, iar $\beta$ este numărul mediu de contacte eficiente pe unitate de
timp; produsul lor dă numărul de infecții secundare produse pe durata infecțiozității.

\begin{teorema}[Pragul epidemic]
\label{thm:prag_epidemic}
Considerăm o populație inițial complet susceptibilă ($S_0 \approx N$). Atunci:
\begin{enumerate}[(i)]
    \item Dacă $\Rzero > 1$, numărul de infectați crește inițial exponențial
    și are loc o epidemie ($I(t)$ atinge un maxim $I_{\max} > I_0$).
    \item Dacă $\Rzero \leq 1$, $I(t)$ este imediat descrescătoare și nu are
    loc nicio epidemie.
\end{enumerate}
\end{teorema}

\begin{proof}
Din $dI/dt|_{t=0} = I_0(\beta S_0/N - \gamma) \approx I_0(\beta - \gamma)$
(pentru $S_0 \approx N$), concluzia urmează din semnul expresiei $\beta - \gamma$:
\begin{equation}
    \diff{I}{t}\bigg|_{t=0} > 0 \iff \beta > \gamma \iff \frac{\beta}{\gamma} > 1 \iff \Rzero > 1.
    \qedhere
\end{equation}
\end{proof}

\subsection{Pragul de imunizare în turmă (herd immunity)}
\label{subsec:herd_immunity}

\begin{definitie}[Imunitate de grup]
Imunitatea de grup (\emph{herd immunity}) se referă la protecția indirectă a
persoanelor susceptibile, generată atunci când o fracție suficient de mare
din populație este imunizată (fie natural, fie prin vaccinare), astfel încât
transmiterea susținută a agentului patogen devine imposibilă.
\end{definitie}

Epidemia se stinge când numărul de reproducere efectiv devine subunitar:
\begin{equation}
    \mathcal{R}(t) = \Rzero \cdot \frac{S(t)}{N} < 1 \iff \frac{S(t)}{N} < \frac{1}{\Rzero}.
\end{equation}
Dacă o fracție $p$ din populație este imunizată (cu $S/N = 1 - p$), condiția
de blocare a epidemiei devine:
\begin{equation}
    (1 - p) \cdot \Rzero < 1 \iff p > 1 - \frac{1}{\Rzero}.
\end{equation}

\begin{definitie}[Pragul de imunizare în turmă]
Fracția minimă a populației care trebuie imunizată pentru a preveni propagarea
epidemiei este:
\begin{equation}
\label{eq:herd_immunity}
    p_c = 1 - \frac{1}{\Rzero}.
\end{equation}
\end{definitie}

\begin{exemplu}
Pentru rujeolă, $\Rzero \approx 12$--$18$, deci $p_c \approx 92$--$94\%$.
Pentru COVID-19 (tulpina originală), $\Rzero \approx 2$--$3$, deci
$p_c \approx 50$--$67\%$. Variantele ulterioare (Delta, Omicron) au $\Rzero$
mai mare, crescând corespunzător pragul de imunizare.
\end{exemplu}

\subsection{Ecuația dimensiunii finale a epidemiei}
\label{subsec:final_size}

O caracteristică remarcabilă a modelului SIR este că permite calculul analitic al
\emph{dimensiunii finale} a epidemiei --- fracția totală din populație care este
infectată în cursul epidemiei --- fără a rezolva explicit sistemul diferențial.

\begin{teorema}[Ecuația dimensiunii finale]
\label{thm:final_size}
Fie $z = 1 - S(\infty)/N$ fracția din populație infectată cumulativ pe durata
epidemiei, presupunând $S_0 = N$ și $R_0 = 0$. Atunci $z$ este unica rădăcină
pozitivă a ecuației transcendentale:
\begin{equation}
\label{eq:final_size}
    z = 1 - e^{-\Rzero z}.
\end{equation}
\end{teorema}

\begin{proof}
Împărțind ecuația pentru $S$ la ecuația pentru $R$:
\begin{equation}
    \frac{dS}{dR} = \frac{-\beta SI/N}{\gamma I} = -\frac{\beta}{\gamma N} S = -\frac{\Rzero}{N} S.
\end{equation}
Această ecuație separabilă în $S$ și $R$ se integrează direct:
\begin{equation}
    \int_{S_0}^{S(t)} \frac{dS'}{S'} = -\frac{\Rzero}{N}\int_0^{R(t)} dR' \implies
    \ln\frac{S(t)}{S_0} = -\frac{\Rzero}{N} R(t).
\end{equation}
Deci $S(t) = S_0 \exp\!\left(-\Rzero R(t)/N\right)$.

La $t \to \infty$, $I(\infty) = 0$ (demonstrat mai jos), deci $R(\infty) = N - S(\infty)$
(pentru $R_0 = 0$). Substituind cu $S_0 = N$:
\begin{equation}
    S(\infty) = N e^{-\Rzero(N - S(\infty))/N}.
\end{equation}
Notând $z = 1 - S(\infty)/N$ și împărțind prin $N$:
\begin{equation}
    1 - z = e^{-\Rzero z} \implies z = 1 - e^{-\Rzero z}.
\end{equation}

\textbf{Că $I(\infty) = 0$:} Funcția $R(t)$ este monoton crescătoare și mărginită
de $N$, deci $\lim_{t\to\infty}R(t)$ există. Dacă $I(\infty) = \ell > 0$, atunci
$R'(t) = \gamma I(t) \geq \gamma\ell > 0$ pentru $t$ suficient de mare, ceea ce
implică $R(t) \to \infty$, contradicție. Deci $I(\infty) = 0$.

\textbf{Unicitatea rădăcinii pozitive:} Fie $g(z) = 1 - e^{-\Rzero z} - z$.
Avem $g(0) = 0$, $g'(0) = \Rzero - 1 > 0$ (dacă $\Rzero > 1$), și
$g(1) = 1 - e^{-\Rzero} - 1 = -e^{-\Rzero} < 0$. Deoarece $g''(z) = -\Rzero^2
e^{-\Rzero z} < 0$, funcția $g$ este concavă, deci are exact o singură rădăcină
în $(0,1)$.
\end{proof}

\begin{observatie}
Ecuația \eqref{eq:final_size} este \emph{transcendentală} și nu admite soluție
în formă închisă pentru $\Rzero \neq 1$. Se rezolvă numeric, de exemplu prin
metoda Newton-Raphson sau prin iterație de punct fix.
\end{observatie}

\begin{corolar}[Rata de atac finală]
Dacă $\Rzero > 1$, rata de atac finală $z^*$ satisface $z^* \in (0, 1)$ și
crește monoton cu $\Rzero$. Limita $\lim_{\Rzero \to \infty} z^* = 1$.
\end{corolar}

\subsection{Extensii ale modelului SIR}
\label{subsec:sir_extensii}

Modelul SIR de bază poate fi extins în mai multe direcții. Introducerea nașterilor
și mortalității naturale permite studiul echilibrelor endemice pe termen lung.
Structura pe vârste poate fi inclusă prin matrice de contact, iar propagarea
geografică prin modele cu patch-uri sau ecuații diferențiale parțiale. Pentru
populații mici sau pentru fazele incipiente ale unei epidemii, modelele stochastice
devin mai potrivite decât formularea deterministă. Intervențiile precum vaccinarea,
izolarea și tratamentul se pot introduce ca fluxuri suplimentare între
compartimente.

% ============================================================
\section{Metode Numerice pentru Ecuații Diferențiale Ordinare}
\label{sec:metode_numerice}
% ============================================================

\subsection{Clasificarea metodelor numerice}
\label{subsec:clasificare_numerice}

Soluția analitică a sistemului SIR nu există în formă închisă (cu excepția
cazurilor degenerate), ceea ce face indispensabilă rezolvarea numerică.
O metodă numerică pentru problema cu valori inițiale (PVI):
\begin{equation}
\label{eq:pvi}
    \diff{\mathbf{y}}{t} = \mathbf{f}(t, \mathbf{y}), \quad \mathbf{y}(t_0) = \mathbf{y}_0,
\end{equation}
generează o secvență de aproximații $\mathbf{y}_n \approx \mathbf{y}(t_n)$ la
punctele $t_n = t_0 + nh$.

Metodele numerice pot fi privite din mai multe unghiuri. Metodele explicite
calculează direct $\mathbf{y}_{n+1}$ din valorile anterioare, în timp ce metodele
implicite includ necunoscuta și în membrul drept, necesitând rezolvarea unui sistem
la fiecare pas. Metodele un-pas folosesc doar starea curentă, pe când metodele
multi-pas folosesc și stări anterioare. Ordinul de convergență $p$ determină
comportamentul erorii globale, $\|\mathbf{y}(t_n)-\mathbf{y}_n\|=\mathcal{O}(h^p)$.
În sfârșit, metodele cu pas adaptiv ajustează $h$ pe baza unei estimări a erorii
locale, spre deosebire de schemele cu pas fix.

\subsection{Metoda Euler explicită}
\label{subsec:euler}

Metoda Euler (1768) este cea mai simplă metodă numerică un-pas, obținută
prin truncherea seriei Taylor la ordinul întâi:
\begin{equation}
\label{eq:euler}
    \mathbf{y}_{n+1} = \mathbf{y}_n + h \cdot \mathbf{f}(t_n, \mathbf{y}_n).
\end{equation}

\begin{teorema}[Eroarea globală a metodei Euler]
\label{thm:euler_eroare}
Dacă $\mathbf{f}$ este Lipschitz în $\mathbf{y}$ cu constanta $L$ și
$\mathbf{y} \in C^2([t_0, T])$, eroarea globală a metodei Euler satisface:
\begin{equation}
\begin{split}
    \max_{n: t_n \leq T} \|\mathbf{y}(t_n) - \mathbf{y}_n\|
    &\leq \frac{h}{2L} \max_{t \in [t_0, T]} \|\mathbf{y}''(t)\| \\
    &\quad \cdot \left(e^{L(T - t_0)} - 1\right)
    = \mathcal{O}(h).
\end{split}
\end{equation}
\end{teorema}

Eroarea locală de trunchiere (pe un singur pas) este $\mathcal{O}(h^2)$, deci
metoda are ordinul local 2 și ordinul global 1.

Aplicată sistemului SIR cu $\mathbf{y}_n = (S_n, I_n, R_n)^\intercal$,
schema Euler este:

\begin{algorithm}
\caption{Metoda Euler pentru sistemul SIR}
\label{alg:euler_sir}
\begin{lstlisting}[style=pseudocode]
Input:  beta, gamma, N, S0, I0, R0, t0, tf, h
Output: Secventa (tn, Sn, In, Rn) pentru tn in [t0, tf]

(S, I, R) <- (S0, I0, R0);  t <- t0
while t < tf do
    dS <- -beta * S * I / N
    dI <-  beta * S * I / N - gamma * I
    dR <-  gamma * I
    S  <- S + h * dS
    I  <- I + h * dI
    R  <- R + h * dR
    t  <- t + h
    Inregistreaza (t, S, I, R)
end while
\end{lstlisting}
\end{algorithm}

\textbf{Stabilitate.} Analiza stabilității metodei Euler pe ecuația test
$y' = \lambda y$ ($\lambda \in \mathbb{C}$) conduce la factorul de amplificare
$z = 1 + h\lambda$. Metoda este stabilă dacă $|1 + h\lambda| \leq 1$, ceea ce
definește domeniul de stabilitate (disc de rază 1 centrat în $-1$ în planul
$h\lambda$). Această restricție limiteaza pasul de integrare $h$, în special
pentru probleme \emph{stiff}.

\subsection{Metoda Runge-Kutta de ordinul 4 (RK4)}
\label{subsec:rk4}

Metodele Runge-Kutta sunt metode un-pas explicite care realizează mai multe
evaluări ale funcției $\mathbf{f}$ la fiecare pas, atingând ordine superioare
fără a recurge la derivate de ordin înalt ale soluției.

\subsubsection{Derivare și formulă}

Metoda RK4 clasică (cu 4 stadii) este dată de:
\begin{equation}
\label{eq:rk4}
\begin{aligned}
    \mathbf{k}_1 &= h \cdot \mathbf{f}(t_n, \mathbf{y}_n), \\
    \mathbf{k}_2 &= h \cdot \mathbf{f}\!\left(t_n + \tfrac{h}{2},\; \mathbf{y}_n + \tfrac{\mathbf{k}_1}{2}\right), \\
    \mathbf{k}_3 &= h \cdot \mathbf{f}\!\left(t_n + \tfrac{h}{2},\; \mathbf{y}_n + \tfrac{\mathbf{k}_2}{2}\right), \\
    \mathbf{k}_4 &= h \cdot \mathbf{f}(t_n + h,\; \mathbf{y}_n + \mathbf{k}_3), \\
    \mathbf{y}_{n+1} &= \mathbf{y}_n + \frac{1}{6}\left(\mathbf{k}_1 + 2\mathbf{k}_2 + 2\mathbf{k}_3 + \mathbf{k}_4\right).
\end{aligned}
\end{equation}

Coeficienții metodei sunt reprezentați compact prin \emph{tabloul Butcher}:
\begin{equation}
\label{eq:butcher_rk4}
\renewcommand{\arraystretch}{1.4}
\begin{array}{c|cccc}
    0   &     &     &     &   \\
    1/2 & 1/2 &     &     &   \\
    1/2 & 0   & 1/2 &     &   \\
    1   & 0   & 0   & 1   &   \\
    \hline
        & 1/6 & 1/3 & 1/3 & 1/6
\end{array}
\end{equation}

\begin{teorema}[Ordinul metodei RK4]
\label{thm:rk4_ordin}
Metoda \eqref{eq:rk4} are ordinul de convergență $p = 4$: dacă
$\mathbf{f} \in C^4$ și $\mathbf{y} \in C^5([t_0, T])$, atunci eroarea globală
satisface:
\begin{equation}
    \max_{n} \|\mathbf{y}(t_n) - \mathbf{y}_n\| = \mathcal{O}(h^4).
\end{equation}
\end{teorema}

\begin{proof}[Schiță de demonstrație]
Eroarea locală de trunchiere $\tau_n$ se obține substituind soluția exactă
$\mathbf{y}(t_n)$ în schema \eqref{eq:rk4} și expandând în serie Taylor.
Se poate arăta că primele 4 condiții de ordin (condițiile de arbore Butcher
pentru $p = 4$) sunt satisfăcute de coeficienții \eqref{eq:butcher_rk4},
implicând $\tau_n = \mathcal{O}(h^5)$. Eroarea globală este
$\mathcal{O}(h^4)$ prin acumularea $\sim 1/h$ de pași. Detalii complete
în \cite{hairer1993}.
\end{proof}

\begin{algorithm}
\caption{Metoda Runge-Kutta de ordinul 4 pentru sistemul SIR}
\label{alg:rk4_sir}
\begin{lstlisting}[style=pseudocode]
Input:  f, y0, t0, tf, h
Output: Secventa (tn, yn)

y <- y0;  t <- t0
while t < tf do
    k1 <- h * f(t,       y           )
    k2 <- h * f(t + h/2, y + k1/2    )
    k3 <- h * f(t + h/2, y + k2/2    )
    k4 <- h * f(t + h,   y + k3      )
    y  <- y + (k1 + 2*k2 + 2*k3 + k4) / 6
    t  <- t + h
end while
\end{lstlisting}
\end{algorithm}

\subsection{Metode Runge-Kutta adaptive (RK45 / Dormand-Prince)}
\label{subsec:rk_adaptive}

Metodele cu pas fix impun utilizatorului alegerea unui $h$ adecvat, care este
o sarcină dificilă în practică. Metodele \emph{adaptive} ajustează automat
pasul $h$ pe baza unei estimări a erorii locale, asigurând că eroarea rămâne
sub o toleranță prestabilită $\varepsilon$.

\subsubsection{Metode Runge-Kutta încorporate}

O \emph{pereche RK încorporată} de ordine $p$ și $p+1$ calculează doi
estimatori pentru $\mathbf{y}_{n+1}$ folosind \emph{același} set de evaluări:
\begin{equation}
    \hat{\mathbf{y}}_{n+1} = \mathbf{y}_n + h\sum_{i=1}^s b_i \mathbf{k}_i \quad \text{(ordinul } p\text{)},
\end{equation}
\begin{equation}
    \tilde{\mathbf{y}}_{n+1} = \mathbf{y}_n + h\sum_{i=1}^s \tilde{b}_i \mathbf{k}_i \quad \text{(ordinul } p+1\text{)},
\end{equation}
cu $\mathbf{k}_i = \mathbf{f}(t_n + c_i h, \mathbf{y}_n + h\sum_{j<i} a_{ij} \mathbf{k}_j)$.
Estimatorul de eroare locală este:
\begin{equation}
    \mathbf{e}_{n+1} = \tilde{\mathbf{y}}_{n+1} - \hat{\mathbf{y}}_{n+1} = h\sum_{i=1}^s (b_i - \tilde{b}_i)\mathbf{k}_i.
\end{equation}

\subsubsection{Metoda Dormand-Prince (DOPRI5 / RK45)}

Perechea \textbf{Dormand-Prince} \cite{dormandprince1980}, implementată ca
\texttt{ode45} în MATLAB, \texttt{RK45} în SciPy (Python), \texttt{Tsit5}
în Julia și \texttt{lsoda}/\texttt{ode45} în R (pachetul \texttt{deSolve}),
utilizează 6 evaluări de funcție și are ordinul principal 5 cu estimator
de eroare de ordinul 4.

Tabloul Butcher complet al metodei Dormand-Prince este:

\begin{equation}
\label{eq:butcher_dopri5}
\renewcommand{\arraystretch}{1.4}
\resizebox{\columnwidth}{!}{$
\begin{array}{c|ccccccc}
    0       &              &             &            &           &             &          & \\
    1/5     & 1/5          &             &            &           &             &          & \\
    3/10    & 3/40         & 9/40        &            &           &             &          & \\
    4/5     & 44/45        & -56/15      & 32/9       &           &             &          & \\
    8/9     & 19372/6561   & -25360/2187 & 64448/6561 & -212/729  &             &          & \\
    1       & 9017/3168    & -355/33     & 46732/5247 & 49/176    & -5103/18656 &          & \\
    1       & 35/384       & 0           & 500/1113   & 125/192   & -2187/6784  & 11/84    & \\
    \hline
    \hat{b} & 35/384       & 0           & 500/1113   & 125/192   & -2187/6784  & 11/84    & 0 \\
    \tilde{b}& 5179/57600  & 0           & 7571/16695 & 393/640   & -92097/339200 & 187/2100 & 1/40
\end{array}
$}
\end{equation}

\subsubsection{Control adaptiv al pasului}

Pasul $h$ este actualizat după fiecare pas pe baza normei erorii:
\begin{equation}
    \text{err} = \left\|\frac{\mathbf{e}_{n+1}}{\text{atol} + \text{rtol} \cdot |\mathbf{y}_{n+1}|}\right\|_\infty,
\end{equation}
unde $\text{atol}$ și $\text{rtol}$ sunt toleranțe absolute și relative.
Pasul următor este calculat prin formula:
\begin{equation}
    h_{\text{nou}} = h \cdot \min\!\left(f_{\max},\; \max\!\left(f_{\min},\; S \cdot \text{err}^{-1/(p+1)}\right)\right),
\end{equation}
cu factorul de siguranță $S = 0.9$ și limitele $f_{\min} = 0.1$, $f_{\max} = 10$.

\subsection{Stabilitate numerică și probleme stiff}
\label{subsec:stiff}

\begin{definitie}[Problemă stiff]
\label{def:stiff}
O problemă cu valori inițiale \eqref{eq:pvi} este numită \emph{stiff} dacă
valorile proprii ale Jacobianului $\partial \mathbf{f}/\partial \mathbf{y}$ au
componente cu parte reală negativă de magnitudine mare, impunând pași de
integrare $h$ extrem de mici pentru metodele explicite din motive de stabilitate,
chiar dacă soluția variază lent.
\end{definitie}

\textbf{Aplicație la modelul SIR.} Jacobianul sistemului SIR \eqref{eq:sir2d} este:
\begin{equation}
    J(S, I) = \begin{pmatrix}
        -\beta I/N & -\beta S/N \\
        \beta I/N & \beta S/N - \gamma
    \end{pmatrix}.
\end{equation}
Valorile proprii depind de $(S, I)$. La debutul epidemiei ($S \approx N$, $I$ mic),
valorile proprii sunt aproximativ $0$ și $\beta - \gamma$, deci problema nu este
stiff. Modelul SIR de bază nu este stiff pentru parametri tipici epidemiologici.
Totuși, extensiile cu mai multe compartimente și rate foarte diferite pot
produce probleme stiff, necesitând metode implicite (Runge-Kutta implicit, BDF).

\begin{observatie}
Metodele implicite (de exemplu, metoda Euler implicită, Crank-Nicolson, Radau)
au domenii de stabilitate nelimitate și sunt preferate pentru probleme stiff,
în ciuda costului computațional superior (rezolvarea unui sistem nonlinear
la fiecare pas).
\end{observatie}

% ============================================================
\section{Date Epidemiologice — Surse și Preprocesare}
\label{sec:date_epidemiologice}
% ============================================================

\subsection{Surse de date epidemiologice}
\label{subsec:surse_date}

Implementarea și validarea modelelor epidemiologice computaționale necesită date
empirice de calitate. Sursele folosite frecvent includ rapoartele Organizației
Mondiale a Sănătății, seriile de supraveghere publicate de ECDC, setul de date COVID-19
al Johns Hopkins University CSSE și colecțiile armonizate de Our World in Data.
Această lucrare folosește două seturi OWID complementare: setul COVID-19 (sursa
principală, cu acoperire globală pe 2020--2023) și setul mpox (focar 2022--2024,
cu o dinamică complet diferită ca scară, durată infecțioasă și raportare). Ambele
oferă serii temporale publice, comparabile între țări și ușor de prelucrat
într-un format comun.

\subsection{Structura datelor și preprocesare}
\label{subsec:preprocesare}

Datele epidemiologice brute conțin frecvent valori lipsă, raportări retrospective,
efecte de weekend și schimbări administrative ale unităților geografice. Înainte de
utilizarea lor într-un model compartimental, seriile trebuie ordonate temporal,
corectate pentru valori imposibile și aduse la o granularitate compatibilă cu pasul
de integrare ales. În caz contrar, erorile de raportare pot fi confundate cu variații
reale ale dinamicii infecției.

\subsection{Estimarea parametrilor din date reale}
\label{subsec:estimare_parametri}

Parametrii $\beta$ și $\gamma$ ai modelului SIR pot fi estimați din date observate
prin minimizarea erorii pătratice dintre seria raportată și soluția numerică, prin
estimare de verosimilitate maximă cu un model de zgomot adecvat sau prin inferență
bayesiană. Alegerea metodei depinde de natura datelor: pentru incidență zilnică,
modelele Poisson sau binomial negativ sunt mai potrivite decât o eroare gaussiană
simplă, în timp ce inferența bayesiană oferă o descriere explicită a incertitudinii.

Estimatorul $\hat{\gamma}$ poate fi obținut direct dacă durata medie a
perioadei infecțioase este disponibilă din studii clinice:
$\hat{\gamma} = 1 / T_{\text{infecțios}}$. Parametrul $\beta$ este mai
dificil de estimat și depinde de comportamentul social, intervenții și
variante ale agentului patogen.
